\documentclass[a4paper,11pt]{article}
\usepackage[T2A]{fontenc}
\usepackage[utf8]{inputenc}
\usepackage[russian]{babel}
\usepackage{amsmath,amssymb}
\usepackage[margin=19mm,headheight=14pt,footskip=15mm]{geometry}
\usepackage{tabularx,booktabs,array}
\usepackage{xcolor}
\usepackage{tikz}
\usepackage{hyperref}
\usepackage{fancyhdr}
\usepackage{microtype}
\renewcommand{\familydefault}{\sfdefault}
\definecolor{accent}{HTML}{006B72}
\definecolor{muted}{HTML}{444F58}
\hypersetup{colorlinks=true,linkcolor=accent,urlcolor=accent,pdftitle={Ревью домашней работы — 22.11.2026}}
\setlength{\parindent}{0pt}
\setlength{\parskip}{0.22em}
\setlength{\abovedisplayskip}{5pt}
\setlength{\belowdisplayskip}{5pt}
\setlength{\abovedisplayshortskip}{3pt}
\setlength{\belowdisplayshortskip}{3pt}
\renewcommand{\arraystretch}{1.22}
\pagestyle{fancy}
\fancyhf{}
\fancyfoot[C]{\footnotesize\color{muted} Лёвин Артём Александрович эксклюзивно для Володи Хачатуряна}
\fancyfoot[R]{\footnotesize\thepage}
\renewcommand{\headrulewidth}{0pt}
\renewcommand{\footrulewidth}{0pt}
\newcommand{\keyidea}{\textcolor{blue!60!black}{\textbf{Ключевая идея.}}\ }
\newcommand{\errorlabel}{\textcolor{red!75!black}{\textbf{Ошибка оформления.}}\ }
\newcommand{\outcome}{\textcolor{accent}{\textbf{Итог.}}\ }

\begin{document}
\begin{titlepage}
\thispagestyle{empty}
\centering
\vspace*{47mm}
{\Large\bfseries\color{accent} МАТЕМАТИКА\par}
\vspace{16mm}
{\Huge\bfseries Ревью домашней работы\par}
\vspace{8mm}
{\Large Формулы сокращённого умножения\par}
\vspace{25mm}
{\large 22 ноября 2026 г.\par}
\vfill
{\large Лёвин Артём Александрович\par}
\vspace{2mm}
{\normalsize эксклюзивно для Володи Хачатуряна\par}
\vspace{16mm}
\begin{tikzpicture}
\draw[accent!65,line width=1.1pt] (-5.7,0) .. controls (-3.1,0.85) and (-1.7,-0.55) .. (0,0)
 .. controls (1.8,0.55) and (3.2,-0.85) .. (5.7,0);
\end{tikzpicture}
\end{titlepage}

\section*{Сводная таблица замечаний}
\small
\begin{tabularx}{\linewidth}{@{}p{0.9cm}p{2.45cm}X p{2.9cm}p{3.1cm}@{}}
\toprule
\textbf{№} & \textbf{Тема} & \textbf{Суть замечания} & \textbf{Ваш ответ} & \textbf{Требуемый ответ} \\
\midrule
\hyperref[task:7]{7} & Разность квадратов & На видимой части работы не названо применённое тождество. Правый край записи обрезан. & Разложение начато; название отсутствует. & Разность квадратов; $(5t-4k)(5t+4k)$. \\
\addlinespace
\hyperref[task:8]{8} & Квадрат разности & На видимой части работы не названо тождество. Последний член итоговой строки вышел за край снимка. & Формула применена; название отсутствует. & Квадрат разности; $9a^2-12ab+4b^2$. \\
\bottomrule
\end{tabularx}
\normalsize

\vspace{6mm}
\textcolor{muted}{\footnotesize Примечание. Оценка опирается на различимые фрагменты фотографий. Часть окончаний строк в заданиях 7 и 8 не попала в кадр; символы на обрезанном участке не восстанавливались по догадке.}

\clearpage
\thispagestyle{fancy}
\vspace*{0.37\textheight}
\begin{center}
{\Huge\bfseries\color{accent} Разбор ошибок}
\end{center}
\clearpage

\phantomsection\label{task:7}
\section*{Задание № 7. Разность квадратов}
\textbf{Условие.} Разложите на множители выражение
\[
25t^2-16k^2.
\]
Укажите название использованной формулы.

\textbf{Ваше решение.} На фотографии читается правильное представление выражения как разности двух квадратов:
\[
25t^2-16k^2=(5t)^2-(4k)^2.
\]
Далее видно начало разложения на линейные множители. Окончание строки с множителями частично обрезано.

\textbf{Ваш ответ.} Название формулы на доступной части работы не записано. Полную конечную запись множителей достоверно прочитать нельзя.

\errorlabel В задании требовались два результата: само разложение \textbf{и} название формулы. Второй обязательный элемент в видимой записи отсутствует.

\textbf{Где именно возникает неполнота?} Последний однозначно правильный шаг:
\[
25t^2-16k^2=(5t)^2-(4k)^2.
\]
Следующий необходимый шаг после разложения --- словесно назвать использованное тождество: \textbf{разность квадратов}. Именно этого элемента ответа не хватает.

\textbf{Почему это важно?} Формулы сокращённого умножения различаются по строению выражения. Если нужно разложить \textbf{разность квадратов}, следует увидеть два квадрата со знаком минус между ними. Умение назвать формулу помогает правильно распознавать её и в более сложных примерах.

\keyidea Для любых выражений $A$ и $B$ верно тождество
\[
\boxed{A^2-B^2=(A-B)(A+B).}
\]
Оно называется \textbf{формулой разности квадратов}.

\textbf{Исправление от последнего правильного шага.} Вы уже получили два квадрата. Обозначим $A=5t$, $B=4k$. Тогда
\[
(5t)^2-(4k)^2=(5t-4k)(5t+4k).
\]
После полученного выражения необходимо написать название формулы.

\textbf{Полное правильное решение.}
\begin{align*}
25t^2-16k^2
 &= (5t)^2-(4k)^2 \\
 &= (5t-4k)(5t+4k).
\end{align*}
Здесь $25t^2=(5t)^2$, а $16k^2=(4k)^2$. Использована \textbf{формула разности квадратов}.

\textbf{Проверка раскрытием скобок.}
\begin{align*}
(5t-4k)(5t+4k)
 &=25t^2+20tk-20tk-16k^2\\
 &=25t^2-16k^2.
\end{align*}
Смешанные слагаемые сокращаются, поэтому получаем исходное выражение.

\outcome
\[
\boxed{25t^2-16k^2=(5t-4k)(5t+4k).}
\]
\textbf{Название:} разность квадратов.

\subsection*{Задача на закрепление}
Разложите на множители $49m^2-9n^2$. Укажите название формулы.

\clearpage
\phantomsection\label{task:8}
\section*{Задание № 8. Квадрат разности}
\textbf{Условие.} Раскройте скобки:
\[
(3a-2b)^2.
\]
Укажите название использованной формулы.

\textbf{Ваше решение.} На фотографии отчётливо читается переход
\[
(3a-2b)^2=(3a)^2-2\cdot(3a)\cdot(2b)+(2b)^2.
\]
Ниже видно начало упрощения $9a^2-12ab+\ldots$. Окончание этой строки оказалось за краем изображения.

\textbf{Ваш ответ.} Название формулы на видимой части работы не указано. Последний член упрощённой записи на фотографии прочитать невозможно.

\errorlabel Не указан обязательный словесный ответ --- название формулы сокращённого умножения. При этом читаемый алгебраический переход выполнен корректно.

\textbf{Где именно возникает неполнота?} Последний надёжно читаемый правильный шаг --- раскрытие квадрата через три слагаемых:
\[
(3a)^2-2\cdot(3a)\cdot(2b)+(2b)^2.
\]
Следовало завершить упрощение, проверить три слагаемых и подписать название: \textbf{квадрат разности}.

\textbf{Почему это важно?} При возведении разности в квадрат возникает \textbf{три} слагаемых. Среднее равно удвоенному произведению двух выражений со знаком минус. Название формулы помогает запомнить её полную структуру и не потерять средний член.

\keyidea Для любых выражений $A$ и $B$:
\[
\boxed{(A-B)^2=A^2-2AB+B^2.}
\]
Это \textbf{формула квадрата разности}.

\textbf{Исправление от последнего правильного шага.} В полученной Вами формуле отдельно вычислите три части:
\[
(3a)^2=9a^2,
\qquad
2\cdot3a\cdot2b=12ab,
\qquad
(2b)^2=4b^2.
\]
Средний член сохраняет знак минус.

\textbf{Полное правильное решение.}
\begin{align*}
(3a-2b)^2
  &= (3a)^2-2\cdot(3a)\cdot(2b)+(2b)^2\\
  &= 9a^2-12ab+4b^2.
\end{align*}
Использована \textbf{формула квадрата разности}.

\textbf{Проверка подстановкой.} Возьмём $a=2$, $b=1$. Левая часть:
\[
(3\cdot2-2\cdot1)^2=4^2=16.
\]
Правая часть:
\[
9\cdot2^2-12\cdot2\cdot1+4\cdot1^2=36-24+4=16.
\]
Значения совпали. Такая проверка подтверждает вычисления для выбранных чисел; общее равенство следует из формулы квадрата разности.

\outcome
\[
\boxed{(3a-2b)^2=9a^2-12ab+4b^2.}
\]
\textbf{Название:} квадрат разности.

\subsection*{Задача на закрепление}
Раскройте скобки в выражении $(4p-3q)^2$. Укажите название формулы.

\end{document}
